\chapter{تصمیم‌گیری}
\label{ch:conditions}

\begin{goals}
\begin{itemize}
  \item اجرای دستورها فقط وقتی شرطی برقرار است، با \fcode{اگر}
  \item انتخاب میان دو راه با \fcode{وگرنه} و میان چند راه با
        \fcode{وگرنه شرط}
  \item شرط‌های تودرتو و ترکیب شرط‌ها
  \item انتخاب مقدار با کنشگر \code{?:}
  \item انتخاب میان حالت‌های مشخص با \fcode{برگزین}
\end{itemize}
\end{goals}

\section{برنامه‌ای که تصمیم می‌گیرد}

همه‌ی برنامه‌هایی که تا این‌جا نوشتیم یک راه بیشتر نداشتند: از خط اول شروع
می‌کردند و خط به خط تا آخر می‌رفتند. اما برنامه‌های واقعی باید بتوانند
تصمیم بگیرند. «اگر نمره از ده بیشتر بود، بنویس قبول.» «اگر رمز درست بود،
وارد شو.» «اگر موجودی کافی نبود، هشدار بده.» این «اگر»ها، دقیقاً همان
واژه‌ای است که در سلام به کار می‌رود.

\section{اگر}

ساده‌ترین شکل تصمیم‌گیری این است: اگر شرطی برقرار بود، یک یا چند دستور
اجرا شود، وگرنه هیچ:

\program{ساده‌ترین اگر}
\begin{salam}
کارکرد آغازین:
    نمره := ۱۴
    اگر نمره >= ۱۰:
        چاپ "قبول"
    پایان
    چاپ "پایان برنامه"
پایان
\end{salam}

\begin{outputfa}
قبول
پایان برنامه
\end{outputfa}

ساختار \fcode{اگر} همان الگویی است که در \fcode{کارکرد آغازین} دیدید: یک خط که
با دونقطه تمام می‌شود، چند دستور با تورفتگی، و \fcode{پایان}. آنچه بین
\fcode{اگر} و دونقطه می‌آید، \textbf{شرط} است: یک عبارت منطقی که یا درست
است یا نادرست. اگر درست باشد، دستورهای درون بلوک اجرا می‌شوند؛ اگر نادرست
باشد، سلام از روی کل بلوک می‌پرد و از خط بعد از \fcode{پایان} ادامه می‌دهد.

نمره را در برنامه‌ی بالا به \code{۸} تغییر دهید و دوباره اجرا کنید. این بار
فقط «پایان برنامه» چاپ می‌شود، چون شرط برقرار نبود و دستور \fcode{چاپ
"قبول"} اصلاً اجرا نشد.

\begin{note}
شرط \fcode{اگر} می‌تواند هر چیزی باشد که یک مقدار منطقی می‌دهد: یک مقایسه
مثل \fcode{نمره \lr{>=} ۱۰}، یک متغیر منطقی مثل \fcode{بلیت دارد}، یا
ترکیبی از این‌ها با \fcode{و} و \fcode{یا}. عدد به تنهایی شرط نیست؛ نوشتن
\fcode{اگر نمره:} خطاست.
\end{note}

\section{اگر و وگرنه}

بسیاری از تصمیم‌ها دو راه دارند: یا این، یا آن. واژه‌ی \kwd{وگرنه} راه دوم
را مشخص می‌کند:

\program{اگر و وگرنه}
\begin{salam}
کارکرد آغازین:
    نمره := ۸
    اگر نمره >= ۱۰:
        چاپ "قبول"
    وگرنه:
        چاپ "مردود"
    پایان
پایان
\end{salam}

\begin{outputfa}
مردود
\end{outputfa}

دقت کنید که \fcode{وگرنه} خودش دونقطه دارد، اما \fcode{پایان} جداگانه ندارد؛
یک \fcode{پایان} در پایان، کل ساختار «اگر، وگرنه» را می‌بندد. از دو بلوک،
همیشه دقیقاً یکی اجرا می‌شود: نه هر دو، نه هیچ‌کدام.

\program{زوج یا فرد}
\begin{salam}
کارکرد آغازین:
    عدد := ۲۷
    اگر عدد % ۲ == ۰:
        چاپ عدد, "زوج است"
    وگرنه:
        چاپ عدد, "فرد است"
    پایان
پایان
\end{salam}

\begin{outputfa}
27 فرد است
\end{outputfa}

عدد زوج آن است که باقی‌مانده‌ی تقسیمش بر ۲ صفر باشد. این یکی از
پرکاربردترین شرط‌ها در برنامه‌نویسی است و بارها آن را خواهید دید.

\section{بیش از دو راه: وگرنه شرط}

گاهی راه‌ها بیش از دو تاست. مثلاً می‌خواهیم نمره را به «عالی»، «خوب»،
«قابل قبول» و «مردود» تقسیم کنیم. می‌توان بعد از \fcode{وگرنه} یک شرط تازه
نوشت:

\program{زنجیره‌ی شرط‌ها}
\begin{salam}
کارکرد آغازین:
    نمره := ۱۷
    اگر نمره >= ۱۸:
        چاپ "عالی"
    وگرنه نمره >= ۱۵:
        چاپ "خوب"
    وگرنه نمره >= ۱۰:
        چاپ "قابل قبول"
    وگرنه:
        چاپ "مردود"
    پایان
پایان
\end{salam}

\begin{outputfa}
خوب
\end{outputfa}

سلام شرط‌ها را از بالا به پایین می‌سنجد و به \emph{اولین} شرطی که درست بود
بسنده می‌کند؛ بقیه اصلاً بررسی نمی‌شوند. نمره‌ی ۱۷ از ۱۸ کمتر است، پس شرط
اول نادرست است؛ اما از ۱۵ بیشتر است، پس «خوب» چاپ می‌شود و سلام از بقیه‌ی
شرط‌ها می‌گذرد. \fcode{وگرنه}‌ی آخر که شرط ندارد، وقتی اجرا می‌شود که
هیچ‌کدام از شرط‌ها برقرار نباشند. نوشتن آن اختیاری است.

\begin{warn}
ترتیب شرط‌ها در زنجیره مهم است. اگر برنامه‌ی بالا را با شرط
\fcode{نمره \lr{>=} ۱۰} شروع می‌کردیم، نمره‌ی ۱۷ هم «قابل قبول» می‌شد و
هیچ‌وقت به شرط‌های بعدی نمی‌رسید. شرط‌ها را از خاص‌ترین به عام‌ترین
بنویسید.
\end{warn}

\section{اگر در یک خط}

وقتی بلوک \fcode{اگر} فقط یک دستور کوتاه دارد، می‌توان همه را در یک خط
نوشت. \fcode{پایان} همچنان لازم است:

\program{اگر تک‌خطی}
\begin{salam}
کارکرد آغازین:
    دما := ۳۵
    اگر دما > ۳۰: چاپ "هوا گرم است" پایان
    اگر دما < ۱۰: چاپ "هوا سرد است" پایان
    اگر دما >= ۱۰ و دما <= ۳۰: چاپ "هوا معتدل است" پایان
پایان
\end{salam}

\begin{outputfa}
هوا گرم است
\end{outputfa}

این شکل برای شرط‌های کوتاه و ساده مناسب است. به محض این که بلوک بیش از
یک دستور داشت یا شرط طولانی شد، به شکل چندخطی برگردید.

\section{شرط درون شرط}

درون بلوک \fcode{اگر} می‌توان هر دستوری نوشت، از جمله یک \fcode{اگر} دیگر.
به این کار \textbf{تودرتو} کردن می‌گویند:

\program{شرط تودرتو}
\begin{salam}
کارکرد آغازین:
    سن := ۲۰
    بلیت دارد := درست
    اگر سن >= ۱۸:
        اگر بلیت دارد:
            چاپ "خوش آمدید"
        وگرنه:
            چاپ "لطفاً بلیت تهیه کنید"
        پایان
    وگرنه:
        چاپ "ورود برای زیر ۱۸ سال ممکن نیست"
    پایان
پایان
\end{salam}

\begin{outputfa}
خوش آمدید
\end{outputfa}

به تورفتگی‌ها دقت کنید: \fcode{اگر} درونی یک پله جلوتر است و \fcode{پایان}
هر بلوک، دقیقاً زیر \fcode{اگر} خودش قرار دارد. بدون تورفتگی درست، فهمیدن
این که کدام \fcode{پایان} کدام بلوک را می‌بندد، تقریباً ناممکن می‌شود.

بسیاری از شرط‌های تودرتو را می‌توان با \fcode{و} و \fcode{یا} صاف کرد.
همان برنامه، بی تودرتویی:

\begin{salam}
کارکرد آغازین:
    سن := ۲۰
    بلیت دارد := درست
    اگر سن >= ۱۸ و بلیت دارد:
        چاپ "خوش آمدید"
    وگرنه سن < ۱۸:
        چاپ "ورود برای زیر ۱۸ سال ممکن نیست"
    وگرنه:
        چاپ "لطفاً بلیت تهیه کنید"
    پایان
پایان
\end{salam}

\begin{outputfa}
خوش آمدید
\end{outputfa}

هیچ‌کدام از دو شکل «درست‌تر» نیست. قاعده‌ی سرانگشتی: هر کدام را که با یک
بار خواندن بهتر می‌فهمید، همان را بنویسید.

\section{بزرگ‌ترین از میان چند عدد}

یک الگوی رایج که با \fcode{اگر} و یک متغیر تغییرپذیر ساخته می‌شود: پیدا
کردن بزرگ‌ترین مقدار. اول فرض می‌کنیم اولی بزرگ‌ترین است، بعد بقیه را با آن
می‌سنجیم:

\program{بزرگ‌ترین سه عدد}
\begin{salam}
کارکرد آغازین:
    الف := ۱۲
    ب := ۲۵
    ج := ۱۹
    گذرا بزرگترین := الف
    اگر ب > بزرگترین: بزرگترین = ب پایان
    اگر ج > بزرگترین: بزرگترین = ج پایان
    چاپ "بزرگ‌ترین:", بزرگترین
پایان
\end{salam}

\begin{outputfa}
بزرگ‌ترین: 25
\end{outputfa}

با هر مقداری برای \fcode{الف}، \fcode{ب} و \fcode{ج} امتحان کنید. این
الگو در فصل \ref{ch:loops} و \ref{ch:arrays}، وقتی به جای سه عدد صدها عدد
داشته باشیم، دوباره به کارمان می‌آید.

\section{انتخاب مقدار: کنشگر ؟}

گاهی نمی‌خواهیم دستوری اجرا شود، بلکه می‌خواهیم بر پایه‌ی شرطی، یکی از دو
\emph{مقدار} را انتخاب کنیم. در فصل پیش کنشگر \fcode{شرط ؟ الف : ب} را
دیدید که دقیقاً همین کار را می‌کند:

\program{کنشگر انتخاب}
\begin{salam}
کارکرد آغازین:
    سن := ۱۵
    قیمت بلیت := سن < ۱۲ ؟ ۵۰۰۰۰ : ۱۲۰۰۰۰
    چاپ "قیمت بلیت:", قیمت بلیت
    چاپ سن >= ۱۸ ؟ "بزرگسال" : "نوجوان"
پایان
\end{salam}

\begin{outputfa}
قیمت بلیت: 120000
نوجوان
\end{outputfa}

خط \fcode{قیمت بلیت \lr{:=} سن \lr{<} ۱۲ ؟ ۵۰۰۰۰ : ۱۲۰۰۰۰} را بخوانید:
«قیمت بلیت برابر است با: اگر سن کمتر از ۱۲ است، ۵۰ هزار، وگرنه ۱۲۰ هزار».
همین کار را می‌شد با یک \fcode{اگر} و \fcode{وگرنه} و یک متغیر تغییرپذیر
کرد، اما این شکل کوتاه‌تر است. علامت پرسش می‌تواند فارسی (\fcode{؟}) یا
انگلیسی (\code{?}) باشد.

\begin{tip}
کنشگر \code{?:} را فقط برای انتخاب‌های ساده به کار ببرید. به محض این که
دو یا سه تای آن‌ها تودرتو شدند، به \fcode{اگر} برگردید؛ خوانایی مهم‌تر از
کوتاهی است.
\end{tip}

\section{برگزین: انتخاب میان حالت‌های مشخص}

وقتی یک مقدار را با چند حالت \emph{مشخص} مقایسه می‌کنیم (شماره‌ی روز
هفته، شماره‌ی ماه، کد یک گزینه)، زنجیره‌ی \fcode{اگر} و \fcode{وگرنه}
طولانی و یکنواخت می‌شود. سلام برای این کار واژه‌ی \kwd{برگزین} را دارد:

\program{نام روز هفته}
\begin{salam}
کارکرد آغازین:
    روز := ۵
    برگزین روز:
        ۱: چاپ "شنبه" پایان
        ۲: چاپ "یکشنبه" پایان
        ۳: چاپ "دوشنبه" پایان
        ۴: چاپ "سه‌شنبه" پایان
        ۵: چاپ "چهارشنبه" پایان
        ۶: چاپ "پنجشنبه" پایان
        ۷: چاپ "جمعه" پایان
        وگرنه: چاپ "شماره‌ی روز نادرست است" پایان
    پایان
پایان
\end{salam}

\begin{outputfa}
چهارشنبه
\end{outputfa}

\fcode{برگزین} مقدار \fcode{روز} را با هر یک از حالت‌ها مقایسه می‌کند و
بلوک همان حالتی را که برابر بود اجرا می‌کند. هر حالت یک بلوک کوچک است که
با دونقطه شروع و با \fcode{پایان} بسته می‌شود؛ می‌توان آن را در یک خط یا
چند خط نوشت. \fcode{وگرنه} برای وقتی است که هیچ حالتی برابر نباشد.

\fcode{برگزین} شکل دومی هم دارد که به جای اجرای دستور، یک \emph{مقدار}
می‌دهد. در این شکل به جای دونقطه از \code{=>} استفاده می‌شود و می‌توان
چند حالت را با ویرگول در یک خط آورد:

\program{فصل هر ماه}
\begin{salam}
کارکرد آغازین:
    ماه := ۴
    فصل := برگزین ماه:
        ۱, ۲, ۳ => "بهار"
        ۴, ۵, ۶ => "تابستان"
        ۷, ۸, ۹ => "پاییز"
        ۱۰, ۱۱, ۱۲ => "زمستان"
        وگرنه => "نامعتبر"
    پایان
    چاپ "ماه", ماه, "در فصل", فصل, "است"
پایان
\end{salam}

\begin{outputfa}
ماه 4 در فصل تابستان است
\end{outputfa}

این شکل برای «جدول‌های تبدیل» (عدد به نام، کد به پیام) بسیار خوش‌دست است.
اگر چنین جدولی را با \fcode{اگر} می‌نوشتیم، دوازده شرط و دوازده
\fcode{چاپ} لازم بود.

\section{یک برنامه‌ی کامل‌تر}

شاخص توده‌ی بدنی از تقسیم وزن (کیلوگرم) بر مجذور قد (متر) به دست می‌آید.
این برنامه شاخص را حساب و دسته‌بندی می‌کند:

\program{شاخص توده‌ی بدنی}
\begin{salam}
کارکرد آغازین:
    وزن := ۷۰.۰
    قد := ۱.۷۵
    شاخص := وزن / (قد * قد)
    چاپ "شاخص توده‌ی بدنی:", شاخص

    اگر شاخص < ۱۸.۵:
        چاپ "وضعیت: کم‌وزن"
    وگرنه شاخص < ۲۵.۰:
        چاپ "وضعیت: طبیعی"
    وگرنه شاخص < ۳۰.۰:
        چاپ "وضعیت: اضافه‌وزن"
    وگرنه:
        چاپ "وضعیت: چاق"
    پایان
پایان
\end{salam}

\begin{outputfa}
شاخص توده‌ی بدنی: 22.857142857142858
وضعیت: طبیعی
\end{outputfa}

به این نکته دقت کنید که در زنجیره‌ی شرط‌ها لازم نبود بنویسیم
\fcode{شاخص \lr{>=} ۱۸.۵ و شاخص \lr{<} ۲۵.۰}؛ چون اگر به شرط دوم
رسیده‌ایم، یعنی شرط اول نادرست بوده و شاخص حتماً از ۱۸٫۵ کمتر نیست.
زنجیره، خودش این را تضمین می‌کند.

\begin{summary}
\begin{itemize}
  \item \fcode{اگر شرط:} بلوکی را فقط وقتی اجرا می‌کند که شرط درست باشد.
        \fcode{وگرنه:} راه دوم را می‌دهد و \fcode{وگرنه شرط:} راه‌های
        بیشتر. یک \fcode{پایان} کل ساختار را می‌بندد.
  \item در زنجیره، اولین شرط درست برنده است؛ ترتیب مهم است.
  \item بلوک‌های کوتاه را می‌توان در یک خط نوشت، با \fcode{پایان} در پایان.
  \item \fcode{اگر} می‌تواند درون \fcode{اگر} باشد؛ تورفتگی را رعایت کنید.
        بسیاری از تودرتویی‌ها با \fcode{و} و \fcode{یا} ساده می‌شوند.
  \item \fcode{شرط ؟ الف : ب} یکی از دو مقدار را انتخاب می‌کند.
  \item \fcode{برگزین} برای مقایسه‌ی یک مقدار با چند حالت مشخص است، به شکل
        دستور (با دونقطه) یا مقدار (با \code{=>}).
\end{itemize}
\end{summary}

\section*{تمرین‌ها}

\exercise عددی را در متغیری بگذارید و چاپ کنید که مثبت، منفی یا صفر است.

\exercise سه ضلع یک مثلث را در سه متغیر بگذارید و چاپ کنید که مثلث
متساوی‌الاضلاع (هر سه ضلع برابر)، متساوی‌الساقین (دو ضلع برابر) یا مختلف‌الاضلاع
است.

\exercise ساعت روز (۰ تا ۲۳) را در متغیری بگذارید و پیام مناسب چاپ کنید:
پیش از ۱۲ «صبح بخیر»، از ۱۲ تا ۱۷ «ظهر بخیر»، از ۱۷ تا ۲۰ «عصر بخیر» و
پس از آن «شب بخیر».

\exercise با \fcode{برگزین}، شماره‌ی یک ماه را به تعداد روزهایش تبدیل کنید:
شش ماه اول ۳۱ روز، پنج ماه بعد ۳۰ روز و اسفند ۲۹ روز.

\exercise وزن یک بسته‌ی پستی را در متغیری بگذارید و هزینه‌ی ارسال را حساب
کنید: تا ۱ کیلوگرم ۳۰ هزار تومان، تا ۵ کیلوگرم ۵۰ هزار تومان و بیش از آن،
به ازای هر کیلوگرم ۱۵ هزار تومان. اگر بسته به شهر دیگری می‌رود (یک متغیر
منطقی)، ۲۰ درصد به هزینه اضافه شود.

\exercise برنامه‌ی زیر چه چاپ می‌کند؟ اول حدس بزنید، بعد اجرا کنید و اگر
حدس شما اشتباه بود، دلیلش را پیدا کنید:
\begin{salam}
کارکرد آغازین:
    عدد := ۱۵
    اگر عدد > ۱۰:
        چاپ "بزرگ‌تر از ده"
    وگرنه عدد > ۵:
        چاپ "بزرگ‌تر از پنج"
    پایان
    اگر عدد % ۵ == ۰: چاپ "مضرب پنج" پایان
    اگر عدد % ۳ == ۰: چاپ "مضرب سه" پایان
پایان
\end{salam}
